\section{Conclusion}
We investigated a latent forecasting framework for longitudinal pCR prediction from breast MRI using matched direct and forecast-based controls. In the retrospectively selected seed44 analysis, forecast-attention readout produced the highest mean AUROC and AP among the principal variants, while auxiliary-supervised direct readout retained the highest mean F1. The complete model's calibration was worse, and paired patient-bootstrap intervals included zero for the principal differences. The findings motivate further study of future-state supervision and forecast-based readout, with prospective model selection, stronger baselines, complete upstream provenance, and independent validation required before a confirmatory or clinical-use claim.
